% Notes of the Montreal seminar, 1970 (cote 129, typed part): archive cote 129, site page 21 (PDF page 22).
% typescript. Draft compared: none.
% AI-transcribed by Claude Opus 5.5 (claude-opus-5-5, adaptive thinking, effort high), Message Batches API, batch mtl129_full (msgbatch_01AGnDKHwy6ATCM7ehUa6263), 2026-11-13; not yet checked by eye.
\scanpage{129/21}
Notons qu'en caractéristique $p$ , les morphismes de transition $TF$ ne
sont autres que \uline{la multiplication par} $p$ puisque l'on a $TFR = FTR = FV = p.Id$ .

\uline{Remarque}. Si $A$ est un anneau \uline{parfait} de caractéristique $p$ , $u$ in-%
duit un isomorphisme
\[ u(A) : \mathbb{W}_{\underset{\to}{\mathbb{F}}_{p}}(A) \to \mathbb{W}'_{\underset{\to}{\mathbb{F}}_{p}}(A) \]
puisque c'est vrai pour chaque $F^{n}(A)$ .

4. \uline{L'anneau} $W$

4.1. On rappelle que (voir J.-P. Serre ``Corps locaux'' II, §6, th.7)
que si $k$ est un corps parfait de caractéristique $p$ , l'anneau
\[ W = \mathbb{W}(k) = \underset{\leftarrow}{\lim}\, W_{n} \qquad (Wn = \mathbb{W}_{n}(k)) \]
est un anneau de valuation discrète, complet, de corps résiduel $k$ , d'uni-%
formisante $p = (0,1,0,0,\ldots)$ , dont l'homomorphisme d'augmentation est
donné par la projection canonique
\[ W \to W_{1} = k . \]

On peut d'ailleurs montrer (loc. cit. II, §5, th. 3) que ces propriétés
caractérisent l'anneau $W$ à un isomorphisme unique près (théorème de Co-%
hen).

De plus, le corps des fractions de $W$ ,
\[ K = W[1/p] \]
est de caractéristique nulle.

\mspage{\hand{21}}
